\section{Introduction and Motivation}
\label{sec:intro}

Most work on visual perception for autonomous driving targets paved urban
streets, where lane markings, curbs, and signs give a model strong geometric
and semantic structure~\cite{cordts2016cityscapes}. Off-road terrain offers
none of this. Boundaries between grass, mulch, and gravel are irregular and
often ambiguous, lighting changes abruptly under tree cover, and a camera on a
small vehicle shakes and blurs as it crosses rough
ground~\cite{wigness2019rugd}. Yet many high-value uses of autonomy happen
off-road: search-and-rescue and disaster-relief robots, agricultural and
forestry vehicles, planetary rovers, and military logistics. For these
systems, distinguishing drivable dirt from standing water or a fallen log is a
central perception problem, and a mistake can leave a vehicle stuck or damaged.
The same capability could also support driver-assist cameras on consumer
off-road vehicles that warn about water crossings, loose rock, or obstacles
before the driver commits to a line.

We treat terrain understanding as \emph{semantic segmentation}: assign each
pixel of a forward-facing camera image one of the 24 terrain and object
classes of the RUGD dataset~\cite{wigness2019rugd}. Two properties make this
harder than urban segmentation. First, the label distribution is extremely
skewed: in the RUGD training split, \emph{tree} and \emph{grass} cover 63\% of non-void
labeled pixels, while \emph{water}, \emph{rock}, and \emph{log} together
cover under 0.8\% (Fig.~\ref{fig:classfreq}), even though these rare classes
are the ones most relevant to safe driving. Second, RUGD is small (7,436
annotated frames from 18 video sequences) and its frames are strongly
correlated in time, so careless splitting would leak near-duplicate frames
between training and test.

This project asks three questions that transfer beyond off-road driving:
(i)~how much do the decoder design and the encoder depth matter when both
start from ImageNet-pretrained weights; (ii)~which class-imbalance loss most
improves accuracy on rare, safety-relevant classes without sacrificing overall
accuracy; and (iii)~how much does ImageNet pretraining help on scenes far from
its training distribution. As a stretch goal, we test how a RUGD-trained model
transfers to RELLIS-3D~\cite{jiang2021rellis}, which was recorded with a
different vehicle in a different environment.

Our contributions are:
\begin{itemize}
  \item A controlled comparison of U-Net~\cite{ronneberger2015unet} and
        DeepLabV3+~\cite{chen2018deeplabv3plus} with ImageNet-pretrained
        ResNet-18/50 encoders~\cite{he2016resnet} on RUGD, changing one factor
        per experiment and selecting models on the validation split only. We
        find that \todo{architecture/encoder finding}.
  \item A comparison of four imbalance-aware losses (median-frequency weighted
        cross-entropy, Dice, focal, and cross-entropy plus Dice) against plain
        cross-entropy, reported per class with emphasis on rare,
        safety-relevant terrain. \todo{loss finding, e.g., which loss improves
        water/log/rock IoU and at what cost to mIoU}.
  \item An ablation of ImageNet pretraining and a multi-seed estimate of
        run-to-run variance for the two strongest configurations.
        \todo{pretraining gap and seed std}.
  \item An error analysis built on the normalized confusion matrix, worst-frame
        galleries, and a six-group traversability evaluation, together with a
        documented account of the data-integrity issues we found in RUGD, such
        as a class that appears only in the test split and a validation split
        that lacks four classes.
        \todo{RELLIS-3D cross-dataset result, if completed.}
  \item An open, configuration-driven pipeline in which every experiment is a
        YAML file and every model-selection decision is made automatically
        from validation scores and logged.
\end{itemize}
